%% kapa: exam=1 round=22 year=2011 period=2 form=A subject=accounting subject_ko=회계학 seq=16 number=16
%% source: Q-Net 감정평가사 기출문제 (cst00309) artlSeq=5036790 file=extracted/1008/5036790/2011년 제22회 감정평가사 1차시험 기출문제/2교시.hwp
%% conversion: kordoc 4.24.51 -> tools/gen1.py (auto); review=proofread-sonnet; flags=none
\begin{question}{16}
(주)서울의 20×1년과 20×2년 결산 마감 후 매출원가는 다음과 같다.

{\small\centering
\begin{tabular}{|c|c|c|}
\hline
 & 20×1년 & 20×2년 \\
\hline
매출원가 & ￦1,000,000 & ￦1,500,000 \\
\hline
\end{tabular}\par}

20×3년에 (주)서울의 회계담당자는 20×0년 말 재고자산이 ￦200,000 과소계상 되었고, 20×1년 말 재고자산이 ￦100,000 과소계상 되었음을 알게 되었으며, 동 오류에 대해서 어떠한 수정도 없었음을 확인하였다. 동 오류가 재무제표에 미치는 영향으로 옳지 않은 것은? \par
\begin{choices}
  \item (주)서울의 20×1년 말 재무상태표상 이익잉여금은 ￦100,000만큼 과소계상되어 있다.
  \item (주)서울의 20×2년 포괄손익계산서상 당기순이익이 ￦3,000,000이었다면, 20×2년의 정확한 당기순이익은 ￦2,900,000이다.
  \item (주)서울의 20×1년의 오류수정 후 매출원가는 ￦900,000이다.
  \item (주)서울의 20×2년 말 재무상태표상 이익잉여금은 적정하게 계상되어 있다.
  \item (주)서울의 20×2년의 오류수정 후 매출원가는 ￦1,600,000이다.
\end{choices}
\end{question}
\answer{3}
